\begin{proof}[Proof of \cref{obj:thm:sharp-high-resolution-constant}]
By \cref{obj:ass:overlap}, \(\mathcal E\) is a nondegenerate compact interval and \(e,1-e\geq\varepsilon>0\) throughout \(\mathcal E\). Continuity and strict positivity give
\[
0<\min_{e\in\mathcal E}f(e)
\leq\max_{e\in\mathcal E}f(e)<\infty.
\]
% lean: CausalSmith.PartialID.UnlinkedPropensityAte.exists_boundedScoreDensity_of_continuous

For \(x,z\in\mathcal E\), define the two distortion coefficients
\[
\beta_1(x,z)=\frac{f(x)}{z},
\qquad
\beta_0(x,z)=\frac{f(x)}{1-z}.
\]
Both are continuous and strictly positive on \(\mathcal E^2\). Their diagonal weight is
\[
w(x):=\beta_0(x,x)+\beta_1(x,x)
=\frac{f(x)}{x(1-x)}.
\]
Thus \cref{obj:lem:paired-high-resolution-quantization} applies with two components. In particular, if \(V_K\) denotes its optimal paired cost for these coefficients, then
\[
K V_K\longrightarrow
\frac14\left(\int_{\varepsilon}^{1-\varepsilon}
\sqrt{\frac{f(x)}{x(1-x)}}\,dx\right)^2.
\]
% lean: CausalSmith.PartialID.UnlinkedPropensityAte.pairedQuantization_tendsto
% lean: CausalSmith.PartialID.UnlinkedPropensityAte.highResolutionBeta_diagonal

We identify this paired cost with the release objective. Fix \(K>0\) and a measurable release \(g:\mathcal E\to\{0,\ldots,K-1\}\), and put \(C_r=\{x\in\mathcal E:g(x)=r\}\). For centers \(\mathbf z=(z_{r,a})_{r,a}\in\mathcal E^{\{0,\ldots,K-1\}\times\{0,1\}}\), define
\[
\Phi_K(g,\mathbf z)
:=
\sum_{r=0}^{K-1}\sum_{a=0}^{1}
\int_{C_r}\beta_a(x,z_{r,a})|x-z_{r,a}|\,dx.
\]
The changes of variables \(z_{r,1}=c^{-1}\) and \(z_{r,0}=1-c^{-1}\) each map the coefficient range in \cref{obj:thm:all-label-ambiguity} bijectively onto \(\mathcal E\), and give
\[
f(x)|1-cx|
=\beta_1(x,z_{r,1})|x-z_{r,1}|,
\qquad
f(x)|1-c(1-x)|
=\beta_0(x,z_{r,0})|x-z_{r,0}|.
\]
The cellwise infima are independent. Each has a nonempty set of feasible values bounded below by zero, including when its cell is empty. Hence \cref{obj:thm:all-label-ambiguity} yields
\[
D_H(g)=\inf_{\mathbf z}\Phi_K(g,\mathbf z).
\]
Taking the infimum over releases flattens these independent choices. Their cells are measurable, pairwise disjoint, and cover \(\mathcal E\); conversely, every such \(K\)-cell partition defines a measurable release by assigning its cell indices as labels. Empty cells are allowed in both descriptions. Therefore \cref{obj:def:optimal-k-ambiguity} gives the exact identity
\[
\Delta_K(H)
=\inf_{g\in\mathcal G_K}\inf_{\mathbf z}\Phi_K(g,\mathbf z)
=V_K.
\]
It follows that
\[
K\Delta_K(H)\longrightarrow
\frac14\left(\int_{\varepsilon}^{1-\varepsilon}
\sqrt{\frac{f(x)}{x(1-x)}}\,dx\right)^2.
\]
% lean: CausalSmith.PartialID.UnlinkedPropensityAte.optimalKAmbiguity_eq_optimalPairedCost

To construct the releases, set
\[
A_H:=\int_{\varepsilon}^{1-\varepsilon}\sqrt{w(x)}\,dx.
\]
For each \(K>0\), choose boundaries \(t_{0,K},\ldots,t_{K,K}\) by
\[
t_{0,K}=\varepsilon,\qquad t_{K,K}=1-\varepsilon,\qquad
\int_{\varepsilon}^{t_{j,K}}\sqrt{w(x)}\,dx=\frac{j}{K}A_H
\quad(0\leq j\leq K).
\]
Positivity of \(w\) makes these boundaries strictly increasing. Let \(g_K\) assign label \(j\) to \([t_{j,K},t_{j+1,K})\) for \(j<K-1\), and label \(K-1\) to \([t_{K-1,K},t_{K,K}]\). Put
\[
\zeta_{j,K}:=\frac{t_{j,K}+t_{j+1,K}}2,
\qquad
U_K:=\Phi_K(g_K,\mathbf z)
\quad\text{with }z_{j,0}=z_{j,1}=\zeta_{j,K}.
\]
The construction clause of \cref{obj:lem:paired-high-resolution-quantization} ensures that these cells form a measurable partition, that every midpoint lies in \(\mathcal E\), and that
\[
K U_K\longrightarrow
\frac14 A_H^2.
\]
% lean: CausalSmith.PartialID.UnlinkedPropensityAte.highResolutionOrderedRelease
% lean: CausalSmith.PartialID.UnlinkedPropensityAte.pairedQuantization_tendsto

For every \(K>0\), optimality over releases gives the lower bound, while the chosen centers give a feasible value for each cellwise infimum:
\[
\Delta_K(H)\leq D_H(g_K)
=\inf_{\mathbf z}\Phi_K(g_K,\mathbf z)
\leq U_K.
\]
Multiplication by \(K\geq0\) preserves both inequalities. At \(K=0\), assign zero to the two auxiliary scaled quantities; both bounds then read \(0\leq0\leq0\). The two established limits squeeze \(K D_H(g_K)\) to \(\frac14A_H^2\).
% lean: CausalSmith.PartialID.UnlinkedPropensityAte.optimalKAmbiguity_le_highResolutionOrderedRelease
% lean: CausalSmith.PartialID.UnlinkedPropensityAte.highResolutionOrderedRelease_worstCaseAmbiguity_le_compandingCost
% lean: CausalSmith.PartialID.UnlinkedPropensityAte.tendsto_K_mul_optimalKAmbiguity

Finally, each displayed cell of \(g_K\) is an interval and is measurable. If \(x,z\) carry the same label and \(x\leq y\leq z\), that interval also contains \(y\); thus the label cells have the asserted interval property.
% lean: CausalSmith.PartialID.UnlinkedPropensityAte.highResolutionOrderedRelease_measurable
% lean: CausalSmith.PartialID.UnlinkedPropensityAte.highResolutionOrderedRelease_cell_ordConnected
\end{proof}
